\documentclass[a4paper, serif, 10pt]{moderncv}

%% moderncv
% color and style
\moderncvcolor{black}
\moderncvstyle{banking}

%% page layout/margins
\usepackage[top=2.5cm, bottom=2.5cm, left=1.5cm, right=1.5cm]{geometry}

\nopagenumbers{}

%% Babel config for Norwegian language
% ref:
% - https://latex3.github.io/babel/guides/locale-norwegian.html
\usepackage[norsk]{babel}

%% fontspec
\usepackage{fontspec}

\IfFontExistsTF{Linux Libertine}{
  \setmainfont[Ligatures={TeX, Common}, Numbers=OldStyle]{Linux Libertine}
}{}
\IfFontExistsTF{Linux Libertine O}{
  \setmainfont[Ligatures={TeX, Common}, Numbers=OldStyle]{Linux Libertine O}
}{}

%% CTeX
% CJK support, used to show CJK characters in the resume
%
% - heading=false: disable ctex heading
% - punct=kaiming: use kaiming punctuations styles for CJK
% - scheme=plain: use plain scheme, do not override \normalsize font size
% - space=auto: space settings for CJK characters
%
% CJK fonts are resolved per locale language, e.g. Noto Serif/Sans CJK JP
% for Japanese resumes, with fallback to the next region if not installed
%
% ref:
% - http://ctan.mirrorcatalogs.com/language/chinese/ctex/ctex.pdf
\usepackage[UTF8, heading=false, punct=kaiming, scheme=plain, space=auto]{ctex}

\IfFontExistsTF{Noto Serif CJK SC}{
  \setCJKmainfont{Noto Serif CJK SC}
}{}
\IfFontExistsTF{Noto Sans CJK SC}{
  \setCJKsansfont{Noto Sans CJK SC}
}{}

%% line spacing
\usepackage{setspace}
\setstretch{1.125}

%% URL styling - use normal text instead of monospace
\urlstyle{same}

%% Patch \httplink and \httpslink to handle full URLs
% The original moderncv macros always prepend http:// or https:// to URLs.
% This causes issues when the URL already has a protocol (e.g., https://example.com)
% resulting in malformed URLs like https://https://example.com.
% This patch checks if the URL already contains "://" and uses it directly if so.
% Note: We use \str_set:Nx (x-expansion) to expand macros like \@homepage first.
\ExplSyntaxOn
\str_new:N \g_yamlresume_url_str

\RenewDocumentCommand{\httpslink}{O{}m}{
  \str_set:Nx \g_yamlresume_url_str {#2}
  \str_if_in:NnTF \g_yamlresume_url_str {://}
    { \url{#2} }
    { \url{https://#2} }
}

\RenewDocumentCommand{\httplink}{O{}m}{
  \str_set:Nx \g_yamlresume_url_str {#2}
  \str_if_in:NnTF \g_yamlresume_url_str {://}
    { \url{#2} }
    { \url{http://#2} }
}
\ExplSyntaxOff

\name{Marte Solberg}{}
\title{Operasjonssjef med erfaring fra logistikk, HMS og forbedringsarbeid}
\phone[mobile]{+47 928 45 112}
\email{marte.solberg@example.no}
\homepage{https://www.martesolberg.no}

\address{Storgata 18, Oslo, Oslo, Norge, 0155}{}{}

\extrainfo{{\small \faLinkedin}\ \href{https://www.linkedin.com/in/marte-solberg}{@marte-solberg}}

\begin{document}

\maketitle

\section{Grunnleggende informasjon}

\cvline{}{\begin{itemize}
\item Erfaren operasjonssjef med over 10 års ledererfaring fra logistikk, service og industri.
\item Dokumentert resultat med å øke driftsregularitet, redusere kostnader og bygge velfungerende tverrfaglige team.
\item Sterk i KPI-oppfølging, budsjettansvar, HMS, kontinuerlig forbedring og samarbeid med leverandører og kunder.
\item Strukturert og løsningsorientert leder som trives med å gjøre komplekse driftsprosesser enkle, målbare og skalerbare.
\end{itemize}}
\section{Utdanning}

\cventry{aug. 2014–juni 2016}
        {Master, Strategisk ledelse, Poeng: 4,3/5}
        {Handelshøyskolen BI}
        {\url{https://www.bi.no}}
        {}
        {\begin{itemize}
\item Fordypning i strategisk styring, operativ excellence og endringsledelse.
\item Gjennomførte casebaserte prosjekter med fokus på kostnadsstyring og prosessforbedring.
\item Utviklet analyse- og beslutningsgrunnlag for ledergrupper i praksisnære prosjekter.
\end{itemize}
\textbf{Kurs}: Strategisk økonomistyring, Operasjons- og kvalitetsledelse, Prosjektledelse, Organisasjon og ledelse, Forhandlinger og beslutninger, Bærekraftig verdikjedestyring}

\cventry{aug. 2011–juni 2014}
        {Bachelor, Økonomi og administrasjon, Poeng: 4,1/5}
        {Universitetet i Oslo}
        {\url{https://www.uio.no}}
        {}
        {\begin{itemize}
\item Solid grunnlag i økonomi, analyse og bedriftsadministrasjon.
\item Arbeidet med kvantitative analyser og presentasjon av beslutningsgrunnlag.
\item Deltok aktivt i studentprosjekter om effektiv ressursbruk og organisasjonsutvikling.
\end{itemize}
\textbf{Kurs}: Mikroøkonomi, Bedriftsøkonomi, Statistikk og analyse, Markedsføring, Regnskap og budsjettering}
\section{Arbeidserfaring}

\cventry{jan. 2021–Nå}
        {Operasjonssjef}
        {Nordvik Logistikk AS}
        {\url{https://www.nordviklogistikk.no}}
        {}
        {\begin{itemize}
\item Leder den daglige driften for et terminal- og distribusjonsnett med 85 ansatte og årlig omsetning på 420 MNOK.
\item Ansvar for produksjonsplanlegging, bemanning, HMS, kvalitet og leverandøroppfølging på tvers av fem lokasjoner.
\item Reduserte avvikskostnader med 18 prosent og økte leveringspresisjonen fra 92 til 97 prosent på 18 måneder.
\item Innførte ukentlig driftsdialog med KPI-er for sikkerhet, kvalitet, kostnad og trivsel.
\item Bygget lederteam og innførte kompetanseplan som reduserte gjennomtrekk blant operativt personell med 25 prosent.
\end{itemize}
\textbf{Nøkkelord}: Driftsledelse, Logistikk, HMS, KPI-oppfølging, Bemanningsplanlegging, Kontinuerlig forbedring, Budsjettansvar, Leverandøroppfølging}

\cventry{mars 2017–des. 2020}
        {Driftsleder}
        {Fjordservice Gruppen AS}
        {\url{https://www.fjordservice.no}}
        {}
        {\begin{itemize}
\item Ledet drift av kundeservice og feltoperasjoner med 45 ansatte og ansvar for døgnkontinuerlig bemanning.
\item Utviklet bemanningsmodell som reduserte overtidsbruk med 22 prosent og forbedret responstid med 15 prosent.
\item Drev systematisk HMS-arbeid, gjennomførte risikovurderinger og fulgte opp avvik og læringstiltak.
\item Forhandlet rammeavtaler med underleverandører og oppnådde årlige besparelser på 6,5 MNOK.
\end{itemize}
\textbf{Nøkkelord}: Driftsstyring, Kundeservice, Bemanning, HMS, Risikostyring, Rammeavtaler, Responstid}

\cventry{aug. 2013–feb. 2017}
        {Produksjonsleder}
        {Oslo Industriservice AS}
        {\url{https://www.osloindustriservice.no}}
        {}
        {\begin{itemize}
\item Ledet produksjon og vedlikehold for en fabrikk med 30 ansatte og to skift.
\item Implementerte Lean-prinsipper og standardiserte arbeidsprosesser, noe som økte OEE fra 71 til 82 prosent.
\item Hadde budsjettansvar for drift, vedlikehold og innkjøp på 55 MNOK.
\item Samarbeidet med fagforeninger og verneombud for å sikre god medvirkning og lavt sykefravær.
\end{itemize}
\textbf{Nøkkelord}: Produksjonsledelse, Lean, Vedlikehold, OEE, Budsjett, Arbeidsmiljø}
\section{Språk}

\cvline{Norsk}{Morsmål eller tospråklig nivå \hfill \textbf{Nøkkelord}: Morsmål}
\cvline{Engelsk}{Fullt profesjonelt nivå \hfill \textbf{Nøkkelord}: C1, Forretningsengelsk}
\cvline{Tysk}{Begrenset arbeidskunnskap \hfill \textbf{Nøkkelord}: Grunnleggende}
\section{Ferdigheter}

\cvline{Operasjonsledelse}{Ekspert \hfill \textbf{Nøkkelord}: Driftsstyring, Bemanningsplanlegging, KPI-oppfølging, HMS, Målstyring, Lederutvikling}
\cvline{Logistikk og verdikjede}{Avansert \hfill \textbf{Nøkkelord}: Lagerstyring, Transportplanlegging, Innkjøp, Leverandøroppfølging, Leveringspresisjon, Varelagerstyring}
\cvline{Forbedringsarbeid}{Avansert \hfill \textbf{Nøkkelord}: Lean, Six Sigma, Prosesskartlegging, Rotårsaksanalyse, Standardisering, Kontinuerlig forbedring}
\cvline{Økonomi og styring}{Avansert \hfill \textbf{Nøkkelord}: Budsjettering, Kostnadskontroll, Avviksanalyse, Investeringskalkyler, Rapportering}
\cvline{Prosjektledelse}{Avansert \hfill \textbf{Nøkkelord}: Prosjektstyring, Endringsledelse, Interessentanalyse, Risikostyring, Gevinstrealisering}
\cvline{Systemer og verktøy}{Viderekommen \hfill \textbf{Nøkkelord}: SAP, Power BI, Excel, Microsoft Dynamics 365, Jira, SharePoint}
\section{Utmerkelser}

\cventry{nov. 2023}
        {Nordvik Logistikk AS}
        {Årets operasjonsleder}
        {}
        {}
        {\begin{itemize}
\item Tildelt for dokumenterte resultater innen driftsforbedring, lederutvikling og HMS.
\item Anerkjent for å ha bygget et sterkt tverrfaglig lederteam og økt medarbeidertilfredsheten.
\end{itemize}}

\cventry{mai 2019}
        {Fjordservice Gruppen AS}
        {Forbedringsprisen}
        {}
        {}
        {\begin{itemize}
\item Tildelt for et vellykket forbedringsprosjekt som reduserte responstid og overtidsbruk.
\item Prosjektet ble trukket frem som beste praksis for tverrfaglig driftsutvikling.
\end{itemize}}
\section{Sertifikater}

\cventry{mai 2022}
        {Norsk Lean Senter}
        {Lean Six Sigma Green Belt}
        {\url{https://www.norskleansenter.no}}
        {}
        {}

\cventry{sep. 2021}
        {PeopleCert}
        {PRINCE2 Foundation}
        {\url{https://www.peoplecert.org}}
        {}
        {}

\cventry{feb. 2020}
        {Arbeidsmiljøsenteret}
        {HMS-kurs for ledere}
        {\url{https://www.arbeidsmiljosenteret.no}}
        {}
        {}
\section{Prosjekter}

\cventry{feb. 2022–aug. 2022}
        {Innføring av digitalt bemanningsverktøy for fem terminaler med mål om bedre ressursutnyttelse og lavere overtid.}
        {Digital bemanningsplanlegging}
        {\url{https://www.nordviklogistikk.no/prosjekter/digital-bemanning}}
        {}
        {\begin{itemize}
\item Kartla bemanningsbehov og etablerte felles planleggingsprosess på tvers av lokasjoner.
\item Innførte verktøy og opplærte 40 ledere og planleggere i nye rutiner.
\item Oppnådde 14 prosent lavere overtidsbruk og 9 prosent høyere produktivitet i pilotperioden.
\end{itemize}
\textbf{Nøkkelord}: Bemanningsplanlegging, Digitalisering, Prosessforbedring, Endringsledelse}

\cventry{mars 2018–nov. 2018}
        {Forbedringsprogram for lager og utlasteri med fokus på flyt, standardisering og sikkerhet.}
        {Lean-program for lageravdeling}
        {\url{https://www.fjordservice.no/prosjekter/lean-lager}}
        {}
        {\begin{itemize}
\item Gjennomførte verdistrømsanalyse og identifiserte flaskehalser i vareflyten.
\item Standardiserte arbeidsinstrukser og innførte tavlemøter for daglig styring.
\item Reduserte gjennomsnittlig ordrebehandlingstid med 19 prosent og HMS-avvik med 30 prosent.
\end{itemize}
\textbf{Nøkkelord}: Lean, Verdistrømsanalyse, Standardisering, HMS, Lagerstyring}
\section{Interesser}

\cvline{Friluftsliv}{Fjellturer, Padling, Langrenn, Sykkel}
\cvline{Styrearbeid og mentorering}{Mentorordning, Næringslivsforum, Foreldreutvalg}
\section{Frivillig arbeid}

\cventry{sep. 2019–Nå}
        {Mentor for unge ledere}
        {Oslo Næringsforening}
        {\url{https://www.oslonaring.no}}
        {}
        {\begin{itemize}
\item Mentorerer unge ledere i operativ drift, prioritering og lederutvikling.
\item Bidrar i fagforum for bærekraftig logistikk og kompetanseutvikling.
\end{itemize}}
    

\cventry{apr. 2015–des. 2021}
        {Operativ frivillig}
        {Norsk Folkehjelp}
        {\url{https://www.folkehjelp.no}}
        {}
        {\begin{itemize}
\item Deltok i beredskapsøvelser og støtteoppdrag ved større arrangementer i Oslo.
\item Gjennomførte førstehjelpskurs og bidro til trygg gjennomføring av lokale arrangementer.
\end{itemize}}
    
\end{document}